% !TEX root = ../../main.tex
%% Défi du chapitre 19 — Le problème de Bâle
%% Résultats vérifiés par verifications/defis_16_20.py.

\rubrique{Défi}

\begin{defi}{Le problème de Bâle : $\zeta(2)=\pi^{2}/6$}{2 h}
Posé en 1644 par Pietro Mengoli, résolu par Leonhard Euler en 1735, le
« problème de Bâle » demande la valeur exacte de
$\zeta(2)=\displaystyle\sum_{k=1}^{+\infty}\frac{1}{k^{2}}$. On en donne une
preuve entièrement élémentaire, par la trigonométrie et les relations de
Viète. Soit $m\in\N^{*}$.

\partieExo{A}{Une identité polynomiale}
\begin{questions}
  \item À l'aide de la formule de Moivre appliquée à
    $(\cos\theta+\mathrm{i}\sin\theta)^{2m+1}$, montrer que
    \[
    \sin\bigl((2m+1)\theta\bigr)=\sum_{k=0}^{m}(-1)^{k}\binom{2m+1}{2k+1}
    \cos^{2m-2k}\theta\,\sin^{2k+1}\theta .
    \]
  \item En déduire, pour $\theta\in\intOO{0}{\frac\pi2}$ :
    $\dfrac{\sin\bigl((2m+1)\theta\bigr)}{\sin^{2m+1}\theta}=P_m\bigl(\cot^{2}\theta\bigr)$,
    où $P_m=\sum_{k=0}^{m}(-1)^{k}\binom{2m+1}{2k+1}X^{m-k}$ et
    $\cot=\frac{\cos}{\sin}$.
  \item Montrer que les $m$ réels $x_r=\cot^{2}\frac{r\pi}{2m+1}$
    ($1\leqslant r\leqslant m$) sont des racines distinctes de $P_m$ ; en
    déduire que ce sont toutes ses racines, puis, par les relations de
    Viète :
    \[
    \sum_{r=1}^{m}\cot^{2}\frac{r\pi}{2m+1}=\frac{m(2m-1)}{3} .
    \]
\end{questions}

\partieExo{B}{L'encadrement}
\begin{questions}[resume]
  \item Montrer que $\sin\theta<\theta<\tan\theta$ pour
    $\theta\in\intOO{0}{\frac\pi2}$, et en déduire
    $\cot^{2}\theta<\dfrac{1}{\theta^{2}}<1+\cot^{2}\theta$.
  \item En sommant ces inégalités pour $\theta=\frac{r\pi}{2m+1}$,
    $1\leqslant r\leqslant m$, montrer que
    \[
    \frac{m(2m-1)}{3}<\frac{(2m+1)^{2}}{\pi^{2}}\sum_{r=1}^{m}\frac{1}{r^{2}}
    <\frac{2m(m+1)}{3} .
    \]
  \item En déduire que $\zeta(2)=\dfrac{\pi^{2}}{6}$.
\end{questions}

\medskip
\begin{center}
\begin{tikzpicture}[x=3cm,y=0.42cm,line cap=round]
  \foreach \v in {2,4,6,8}{\draw[otRule,line width=0.3pt] (0,\v) -- (1.62,\v);
    \node[font=\tiny,text=otMuted,anchor=east] at (0,\v) {$\v$};}
  \draw[otMuted,line width=0.5pt,-{Stealth[length=1.6mm]}] (0,0) -- (1.75,0);
  \draw[otMuted,line width=0.5pt,-{Stealth[length=1.6mm]}] (0,0) -- (0,9.2);
  \draw[otMuted,line width=0.4pt,dash pattern=on 2pt off 1.5pt] (1.5708,0) -- (1.5708,8.6);
  \draw[otNavy!45,line width=1pt] plot[domain=0.343:1.5708,samples=120] ({\x},{1/(sin(deg(\x))^2)});
  \draw[otGreenDark,line width=1.2pt] plot[domain=0.343:1.5708,samples=120] ({\x},{1/(\x*\x)});
  \draw[otNavy,line width=1pt] plot[domain=0.36:1.5708,samples=120] ({\x},{(cos(deg(\x))/sin(deg(\x)))^2});
  \node[font=\tiny,text=otMuted,below] at (1.5708,0) {$\frac\pi2$};
  \node[font=\scriptsize,text=otNavy!60,anchor=west] at (0.62,3.6) {$1+\cot^{2}\theta$};
  \node[font=\scriptsize,text=otGreenDark,anchor=west] at (0.86,1.9) {$1/\theta^{2}$};
  \node[font=\scriptsize,text=otNavy,anchor=north east] at (0.78,0.7) {$\cot^{2}\theta$};
\end{tikzpicture}\\[0.5mm]
\legende{Sur $\intOO{0}{\frac\pi2}$, la courbe de $1/\theta^{2}$ (en vert) est
coincée entre celles de $\cot^{2}\theta$ et de $1+\cot^{2}\theta=1/\sin^{2}\theta$.}
\end{center}

\partieExo{C}{Conséquences}
\begin{questions}[resume]
  \item En séparant termes pairs et impairs, en déduire
    $\displaystyle\sum_{k\geqslant0}\frac{1}{(2k+1)^{2}}=\frac{\pi^{2}}{8}$ et
    $\displaystyle\sum_{k\geqslant1}\frac{(-1)^{k+1}}{k^{2}}=\frac{\pi^{2}}{12}$.
  \item On admet l'encadrement du reste
    $\frac{1}{N+1}<\sum_{k>N}\frac{1}{k^{2}}<\frac{1}{N+\frac12}$
    (chapitre 19, exercice 20). Sachant que
    $\sum_{k=1}^{10}\frac1{k^{2}}\approx1{,}549\,768$, en déduire un
    encadrement de $\pi$ d'amplitude inférieure à $5\cdot10^{-3}$.
\end{questions}
\end{defi}
